\chapter{Rozkłady macierzy}
\label{ch:lin15-rozklady-macierzy}
\index{rozkład LU@Rozkład LU}
\index{rozkład Choleskiego@Rozkład Choleskiego}
\index{rozkład QR@Rozkład QR}
\index{rozkład według wartości osobliwych@Rozkład według wartości osobliwych}
\index{rozkład biegunowy@Rozkład biegunowy}

\noindent\textbf{Cel rozdziału.}
W rozdziałach~\ref{ch:lin13-ortogonalnosc-rzuty}
i~\ref{ch:lin14-spektralne} otrzymaliśmy operatory
sprzężone, rzuty ortogonalne oraz twierdzenie
spektralne. Teraz wykorzystamy te wyniki do
rozkładów macierzy stosowanych w rozwiązywaniu
układów liniowych, aproksymacji i teorii operatorów.
Każdy rozkład zostanie odróżniony od innych
przez jawne założenia i zakres jednoznaczności.
Część poświęcona eliminacji Gaussa działa nad
dowolnym ciałem, natomiast rozkłady Choleskiego,
QR, według wartości osobliwych oraz biegunowy
wymagają rzeczywistego lub zespolonego
iloczynu skalarnego.

Przez \(K\) oznaczamy dowolne ciało do chwili,
gdy jawnie ograniczymy je do
\(\mathbb R\) lub \(\mathbb C\).
Macierze mogą mieć rozmiar zerowy; dla
pustych macierzy przyjmujemy standardowe
umowy dotyczące wyznacznika i identyczności.
Dla \(K=\mathbb R,\mathbb C\) symbol \(A^*\)
oznacza sprzężenie transpozycję, a na
\(K^n\) używamy standardowego iloczynu
skalarnego, liniowego w drugim argumencie.

\section{Rozkłady trójkątne i eliminacja Gaussa}
\label{sec:lin15-lu}

\begin{definicja}[Rozkłady LU i PLU]
\label{def:lin15-lu}
Niech \(A\in\operatorname{Mat}_{n\times n}(K)\).
\emph{Rozkładem LU} macierzy \(A\) nazywamy
równość
\[
 A=LU,
\]
gdzie \(L\) jest macierzą dolnotrójkątną
o jedynkach na przekątnej, a \(U\) jest
macierzą górnotrójkątną.
\emph{Rozkładem PLU} nazywamy równość
\[
 PA=LU,
\]
gdzie \(P\) jest macierzą permutacji,
a \(L,U\) spełniają te same warunki.
Jedynki na przekątnej \(L\) służą
ustaleniu normalizacji czynników.
\end{definicja}

\begin{twierdzenie}[Rozkład LU bez zamian wierszy]
\label{thm:lin15-lu}
Niech \(A\in\operatorname{Mat}_{n\times n}(K)\).
Dla \(k=1,\ldots,n\) oznaczmy przez
\(\Delta_k\) wyznacznik podmacierzy
\(A\) złożonej z pierwszych \(k\)
wierszy i pierwszych \(k\) kolumn;
połóżmy \(\Delta_0=1_K\).
Jeżeli
\[
 \Delta_k\ne0_K\qquad(1\le k\le n),
\]
to istnieje dokładnie jeden rozkład
\(A=LU\) z dolnotrójkątną macierzą
\(L\) o jedynkach na przekątnej
i górnotrójkątną macierzą \(U\).
Ponadto
\[
 U_{kk}=\frac{\Delta_k}{\Delta_{k-1}}
 \qquad(1\le k\le n).
\]
Przy \(n=0\) obie macierze są puste,
a teza pozostaje prawdziwa.
\end{twierdzenie}
\begin{proof}
Wymiar \(n=0\) nie wymaga dowodu.
Dla \(n=1\) jedynym rozkładem
jest \(L=(1)\), \(U=A\);
założenie mówi, że
jedyny element \(A\) jest niezerowy.

Załóżmy \(n\ge2\) i zapiszmy
\[
 A=\begin{pmatrix}a&b\\c&C\end{pmatrix},
\]
gdzie \(a\in K\), \(b\) jest
wierszem, \(c\) kolumną,
a \(C\) macierzą kwadratową
stopnia \(n-1\). Z założenia
\(a=\Delta_1\ne0\).
Definiujemy \emph{dopełnienie Schura}
\[
 S=C-ca^{-1}b.
\]
Bezpośrednie mnożenie bloków daje
\[
 A=
 \begin{pmatrix}1&0\\ca^{-1}&I_{n-1}\end{pmatrix}
 \begin{pmatrix}a&b\\0&S\end{pmatrix}.
\]
Dla \(1\le j\le n-1\) taki sam
rozkład blokowy można zastosować
do pierwszych \(j+1\) wierszy
i kolumn \(A\). Własności
wyznacznika macierzy blokowo
trójkątnej dają
\[
 \Delta_{j+1}
 =a\det S_{[1:j,\,1:j]},
\]
gdzie \(S_{[1:j,\,1:j]}\)
jest wiodącą podmacierzą \(S\)
stopnia \(j\). Wobec niezerowości
\(a\) i \(\Delta_{j+1}\)
wszystkie wiodące wyznaczniki
\(S\) są niezerowe.
Z założenia indukcyjnego
istnieje rozkład \(S=L_SU_S\).
Wówczas
\[
 L=\begin{pmatrix}1&0\\ca^{-1}&L_S\end{pmatrix},
 \qquad
 U=\begin{pmatrix}a&b\\0&U_S\end{pmatrix}
\]
spełniają żądane warunki
i \(LU=A\).

Pokażemy jednoznaczność.
W każdym takim rozkładzie
pierwszy wiersz \(U\) jest
pierwszym wierszem \(A\),
ponieważ pierwszy wiersz \(L\)
to \((1,0,\ldots,0)\).
W szczególności \(U_{11}=a\).
Pierwsza kolumna \(L\)
jest wówczas jednoznacznie
wyznaczona przez pierwszą
kolumnę \(A\):
\(L_{i1}=A_{i1}a^{-1}\) dla \(i>1\).
Pozostałe bloki \(L_S,U_S\)
muszą spełniać \(S=L_SU_S\)
i są jednoznaczne przez indukcję.

Wreszcie wiodąca podmacierz
\(A\) stopnia \(k\) jest
iloczynem wiodących podmacierzy
\(L,U\) stopnia \(k\).
Ponieważ pierwsza ma
wyznacznik równy \(1_K\),
otrzymujemy
\(\Delta_k=\prod_{j=1}^{k}U_{jj}\).
Dzieląc tę równość przez
analogiczną dla \(k-1\),
uzyskujemy wzór z tezy.
\end{proof}

\begin{wniosek}[Rozkład LDU]
\label{cor:lin15-ldu}
Jeśli \(A\in\operatorname{Mat}_{n\times n}(K)\)
spełnia założenia
Twierdzenia~\ref{thm:lin15-lu},
to istnieją jednoznaczne macierze
\(L,D,R\) takie, że
\[
 A=LDR,
\]
gdzie \(L\) jest dolnotrójkątna
z jedynkami na przekątnej,
\(D\) diagonalna i odwracalna,
a \(R\) górnotrójkątna
z jedynkami na przekątnej.
Elementy diagonalne \(D\)
są równe
\(D_{kk}=\Delta_k/\Delta_{k-1}\).
\end{wniosek}
\begin{proof}
Niech \(A=LU\) będzie jedynym
rozkładem z poprzedniego twierdzenia.
Wszystkie \(U_{kk}\) są niezerowe.
Kładziemy
\(D=\operatorname{diag}(U_{11},\ldots,U_{nn})\)
oraz \(R=D^{-1}U\).
Wówczas \(R\) jest górnotrójkątna
i ma jedynki na przekątnej,
a \(LDR=LU=A\).
Jeśli \(A=L'D'R'\) jest innym
takim rozkładem, to
\(A=L'(D'R')\) jest rozkładem LU
z tą samą normalizacją.
Jednoznaczność LU daje
\(L'=L\) oraz \(D'R'=U\).
Porównanie przekątnych
wyznacza \(D'=D\), a potem
\(R'=D^{-1}U=R\).
\end{proof}

\begin{przyklad}[Eliminacja bez zamiany wierszy]
\label{prz:lin15-lu}
Dla macierzy
\[
 A=\begin{pmatrix}2&1\\4&3\end{pmatrix}
\]
wiodące wyznaczniki wynoszą
\(\Delta_1=2\) i \(\Delta_2=2\).
Wobec tego rozkład LU
istnieje i jest jednoznaczny:
\[
 L=\begin{pmatrix}1&0\\2&1\end{pmatrix},
 \qquad
 U=\begin{pmatrix}2&1\\0&1\end{pmatrix}.
\]
Równość \(A=LU\)
sprawdzamy przez wymnożenie.
Elementy diagonalne \(U\)
są dokładnie ilorazami
\(\Delta_1/\Delta_0=2\)
i \(\Delta_2/\Delta_1=1\).
\end{przyklad}

\begin{twierdzenie}[Istnienie rozkładu PLU]
\label{thm:lin15-plu}
Dla każdej macierzy kwadratowej
\(A\in\operatorname{Mat}_{n\times n}(K)\)
nad dowolnym ciałem \(K\),
również dla macierzy osobliwej,
istnieje macierz permutacji \(P\),
dolnotrójkątna \(L\) z jedynkami
na przekątnej oraz górnotrójkątna \(U\)
takie, że \(PA=LU\).
Nie twierdzimy, że taki
rozkład jest jednoznaczny.
\end{twierdzenie}
\begin{proof}
Prowadzimy indukcję względem \(n\).
Dla \(n=0\) stosujemy macierze puste.
Niech \(n\ge1\).
Jeżeli pierwsza kolumna \(A\)
jest zerowa, zapiszmy
\[
 A=\begin{pmatrix}0&b\\0&C\end{pmatrix}.
\]
Z indukcji istnieją
\(P',L',U'\) dla \(C\), przy czym
\(P'C=L'U'\).
Połóżmy
\[
 P=\begin{pmatrix}1&0\\0&P'\end{pmatrix},
 \quad
 L=\begin{pmatrix}1&0\\0&L'\end{pmatrix},
 \quad
 U=\begin{pmatrix}0&b\\0&U'\end{pmatrix}.
\]
Wtedy \(PA=LU\),
\(P\) jest permutacją,
a \(L,U\) mają wymagane
kształty.

Jeśli pierwsza kolumna
nie jest zerowa, wybierzmy
permutację wierszy \(P_0\),
która przenosi niezerowy
element tej kolumny
na pierwszą pozycję.
Zapiszmy
\[
 P_0A=\begin{pmatrix}a&b\\c&C\end{pmatrix},
 \qquad a\ne0.
\]
Niech
\(S=C-ca^{-1}b\).
Z założenia indukcyjnego
istnieje
\(P'S=L'U'\).
Definiujemy
\[
 \begin{split}
 P&=\begin{pmatrix}1&0\\0&P'\end{pmatrix}P_0,\\
 L&=\begin{pmatrix}1&0\\P'ca^{-1}&L'\end{pmatrix},
 \qquad
 U=\begin{pmatrix}a&b\\0&U'\end{pmatrix}.
 \end{split}
\]
Macierz \(P\) jest permutacją,
\(L\) jest dolnotrójkątna
z jedynkami na przekątnej,
a \(U\) jest górnotrójkątna.
Bezpośrednie mnożenie daje
\[
 LU=
 \begin{pmatrix}
 a&b\\P'c&P'ca^{-1}b+L'U'
 \end{pmatrix}
 =
 \begin{pmatrix}
 a&b\\P'c&P'C
 \end{pmatrix}
 =PA,
\]
ponieważ
\(L'U'=P'S=P'C-P'ca^{-1}b\).
W obu przypadkach indukcja
jest zamknięta.
\end{proof}

\begin{przyklad}[Konieczność zamiany wierszy]
\label{prz:lin15-pivot}
Niech
\[
 A=\begin{pmatrix}0&2\\1&3\end{pmatrix}.
\]
Nie istnieje jej rozkład LU
z jedynkami na przekątnej \(L\).
Istotnie, z pierwszego wiersza
otrzymalibyśmy
\(U_{11}=0\), \(U_{12}=2\),
a pierwszy element drugiego
wiersza iloczynu \(LU\)
wynosiłby \(L_{21}U_{11}=0\),
chociaż \(A_{21}=1\).
Po zamianie wierszy
\[
 P=\begin{pmatrix}0&1\\1&0\end{pmatrix},
 \qquad
 PA=\begin{pmatrix}1&3\\0&2\end{pmatrix}
 =I_2\begin{pmatrix}1&3\\0&2\end{pmatrix}.
\]
To poprawny rozkład PLU.
\end{przyklad}

\begin{uwaga}[Rozwiązywanie układów trójkątnych]
\label{rem:lin15-triangular}
Dla odwracalnej macierzy \(A\)
o rozkładzie \(PA=LU\)
układ \(Ax=b\) jest równoważny
dwóm układom
\(Ly=Pb\) i \(Ux=y\).
Pierwszy rozwiązujemy
przez podstawianie w przód,
drugi przez podstawianie wstecz.
Nie jest do tego potrzebne
obliczanie macierzy \(A^{-1}\).
W rachunku numerycznym o skończonej
precyzji wybór elementu głównego
ma dodatkowo znaczenie dla stabilności
obliczeń; algebraiczne twierdzenie
o istnieniu PLU samo nie daje
oszacowania błędu zaokrągleń.
\end{uwaga}

\section{Rozkład Choleskiego}
\label{sec:lin15-cholesky}

\begin{twierdzenie}[Rozkład Choleskiego]
\label{thm:lin15-cholesky}
Niech \(K=\mathbb R\) lub \(\mathbb C\),
a \(H\in\operatorname{Mat}_{n\times n}(K)\)
będzie macierzą hermitowską
i dodatnio określoną:
\(H^*=H\) oraz
\(x^*Hx>0\) dla każdego \(x\ne0\).
Istnieje dokładnie jedna
macierz dolnotrójkątna
\(L\in\operatorname{Mat}_{n\times n}(K)\)
o dodatnich rzeczywistych
elementach diagonalnych taka, że
\[
 H=LL^*.
\]
Dla \(n=0\) jest to równość
macierzy pustych.
\end{twierdzenie}
\begin{proof}
Stosujemy indukcję względem \(n\).
Dla \(n=0\) teza jest oczywista.
Dla \(n\ge1\) zapiszmy
\[
 H=\begin{pmatrix}a&b^*\\b&C\end{pmatrix},
\]
gdzie \(a\in\mathbb R\),
\(b\in K^{n-1}\),
a \(C=C^*\).
Dodatnia określoność zastosowana
do pierwszego wektora bazowego
daje \(a>0\).
Tworzymy dopełnienie Schura
\[
 S=C-\frac1a bb^*.
\]
Pokażemy, że \(S\)
jest hermitowska i dodatnio określona.
Hermitowskość wynika z
\(C^*=C\) i \((bb^*)^*=bb^*\).
Dla \(z\in K^{n-1}\setminus\{0\}\)
weźmy
\(x=(-a^{-1}b^*z,z)^{\mathsf T}\).
Wektor \(x\) jest niezerowy,
a mnożenie blokowe daje
\[
 x^*Hx=z^*Sz.
\]
Lewa strona jest dodatnia;
zatem \(z^*Sz>0\).
Na mocy założenia indukcyjnego
istnieje jedyna dolnotrójkątna
macierz \(L_S\) o dodatniej
przekątnej taka, że
\(S=L_SL_S^*\).
Definiujemy
\[
 L=\begin{pmatrix}
 \sqrt a&0\\
 b/\sqrt a&L_S
 \end{pmatrix}.
\]
Wtedy
\[
 LL^*=
 \begin{pmatrix}
 a&b^*\\b&bb^*/a+L_SL_S^*
 \end{pmatrix}
 =H,
\]
a elementy diagonalne \(L\)
są dodatnie.

Dla jednoznaczności rozważmy
dowolną dolnotrójkątną
\(\widetilde L\) z dodatnią
rzeczywistą przekątną
spełniającą
\(H=\widetilde L\widetilde L^*\).
Jej pierwszy element diagonalny
spełnia \(\widetilde L_{11}^2=a\),
więc \(\widetilde L_{11}=\sqrt a\).
Porównanie pierwszej kolumny
poza przekątną daje dokładnie
\(b/\sqrt a\). Pozostały blok
musi zatem spełniać
\(S=\widetilde L_S\widetilde L_S^*\).
Jednoznaczność indukcyjna
daje \(\widetilde L_S=L_S\).
\end{proof}

\begin{przyklad}[Rozkład Choleskiego nad \(\mathbb C\)]
\label{prz:lin15-cholesky}
Niech
\[
 H=\begin{pmatrix}4&2i\\-2i&5\end{pmatrix}.
\]
Wtedy \(H^*=H\), a
\[
 z^*Hz=
 4\left|z_1+\frac i2z_2\right|^2
 +4|z_2|^2>0
 \qquad(z\ne0).
\]
Macierz
\[
 L=\begin{pmatrix}2&0\\-i&2\end{pmatrix}
\]
ma dodatnią przekątną i
\[
 LL^*=
 \begin{pmatrix}2&0\\-i&2\end{pmatrix}
 \begin{pmatrix}2&i\\0&2\end{pmatrix}
 =\begin{pmatrix}4&2i\\-2i&5\end{pmatrix}=H.
\]
Zastosowanie zwykłej transpozycji
zamiast sprzężonej dałoby
w przypadku zespolonym
niepoprawny rozkład.
\end{przyklad}

\begin{uwaga}[Znaczenie dodatniej określoności]
\label{rem:lin15-cholesky-psd}
Dodatnia półokreśloność nie wystarcza
do twierdzenia o jedynym czynniku
trójkątnym z \emph{dodatnią}
przekątną. Na przykład
macierz \(\operatorname{diag}(1,0)\)
jest dodatnio półokreślona,
ale każdy rozkład \(LL^*\)
ma ostatni element diagonalny
równy zeru, nie dodatni.
Ogólne rozkłady dla macierzy
półokreślonych wymagają odrębnego
sformułowania i nie są
używane bez założeń.
\end{uwaga}
